\documentclass[11pt]{article}
\usepackage[margin=1in]{geometry}
\usepackage{amsmath,amssymb,amsthm}
\usepackage{booktabs}
\usepackage{longtable}
\usepackage{array}
\usepackage{enumitem}
\usepackage[T1]{fontenc}
\usepackage{lmodern}
\usepackage{microtype}
\usepackage{xcolor}
\usepackage{hyperref}
\hypersetup{colorlinks=true, linkcolor=blue!50!black, urlcolor=blue!50!black}
\usepackage{listings}
\lstset{basicstyle=\ttfamily\small, columns=fullflexible, keepspaces=true,
  breaklines=true, frame=single, framerule=0.3pt, xleftmargin=4pt, xrightmargin=4pt}
\newcommand{\code}[1]{\texttt{#1}}
\newcommand{\rbar}{\bar r}
\newcommand{\expit}{\operatorname{expit}}
\newcommand{\E}{\mathcal{E}}
\DeclareMathOperator{\logsumexp}{logsumexp}
\setlength{\parskip}{4pt}

\title{The joint amyloid--tau extension:\\ what was added to the code, and what each piece computes}
\author{Implementation note for \texttt{amyloid\_mixture\_cure}}
\date{October 21, 2026}

\begin{document}
\maketitle

\begin{abstract}
\noindent
This note documents, item by item, the code added to the amyloid mixture-cure ODE
package to implement the joint amyloid--tau model proposed in
\code{docs/context/coupling\_and\_copula.pdf}: three lifetime susceptibility
states, a Gaussian copula on the two departure ages, and a tau progression rate
multiplied by the subject's current amyloid burden. For every new function,
model node, setting, data field and test, it states the mathematical object
being computed, the argument conventions, and any deviation from the proposal.
The amyloid-only model (\code{R/00}--\code{R/05}, \code{run\_model.R}) is not
modified; the extension is additive. Nothing here reports a fit: the joint model
has been compiled and smoke-run only.
\end{abstract}

\tableofcontents

\section{Inventory of changes}

\begin{longtable}{@{}p{0.34\linewidth}p{0.62\linewidth}@{}}
\toprule
File & Change \\
\midrule
\endhead
\code{data/raw/dta\_amy.csv}, \code{dta\_tau.csv} & Source CSVs copied into the project. \\
\code{data/build\_joint\_dataset.R} & \textbf{New.} Builds \code{data/adni\_joint.rds} (\S\ref{sec:data}). \\
\code{data/adni\_joint.rds} & \textbf{New.} The dataset the models now load. \\
\code{run\_model.R} & \code{load\_amyloid\_data()} defaults to the joint file; new helper \code{subset\_subjects()}. \\
\code{simulation/sim\_generate.R} & Uses \code{subset\_subjects()} in place of hand-written subsetting. \\
\code{R/00\_settings.R} & \textbf{Added} \code{SETTINGS\_TAU} and \code{SETTINGS\_JOINT} (\S\ref{sec:settings}). \\
\code{R/06\_joint\_densities.R} & \textbf{New.} All joint-model \code{nimbleFunction}s and the two custom distributions (\S\ref{sec:densities}). \\
\code{R/07\_joint\_model.R} & \textbf{New.} \code{build\_joint\_model()}: the NIMBLE model, constants, data, initial values (\S\ref{sec:model}). \\
\code{run\_joint\_model.R} & \textbf{New.} \code{configure\_joint\_samplers()}, \code{fit\_joint\_model()}, \code{joint\_state\_probs()}, \code{joint\_paths()} (\S\ref{sec:runner}). \\
\code{run\_joint\_60k.R} & \textbf{New.} Paste-in script for a 3-chain, 60{,}000-iteration fit. \\
\code{tests/test\_joint.R} & \textbf{New.} 24 deterministic checks (\S\ref{sec:tests}). \\
\code{tests/test\_package.R} & Four checks added for the joint data file and \code{subset\_subjects()}. \\
\code{README.md} & File table updated; new section 10 on the joint model. \\
\code{docs/context/*.pdf} & The three proposal notes and the paper draft, for reference. \\
\bottomrule
\end{longtable}

\section{Notation}

Subjects $i=1,\dots,N$ ($N=1101$); biomarkers $k\in\{A,T\}$; scans
$j=1,\dots,n_{ki}$ at calendar ages $a_{kij}$ with values $y_{kij}$ (SUVR).
Throughout, ``level'' means a biomarker value and ``age'' means calendar age in
years. Quantities inherited from the amyloid model keep their names:
$\mu_{bk}$ (launch level and non-accumulator mean), $c_k$ (positivity threshold),
$r_k(y)=\exp\{B_k(y)^\top\theta_k\}$ (rate curve), $\rbar_k=\max(r_k,r_{\min})$,
$\omega_k(y)=\exp\{V_k(y)^\top\nu_k\}$ (measurement scale), and the clock
\begin{equation}
G_k(y)=\int_{c_k}^{y}\frac{dv}{\rbar_k(v)},\qquad
G_{k,\mathrm{ext}}(x)=\begin{cases}G_k(x), & x\ge\mu_{bk},\\
G_k(\mu_{bk})+(x-\mu_{bk})/\rbar_k(\mu_{bk}), & x<\mu_{bk}.\end{cases}
\label{eq:clock}
\end{equation}
Each subject carries location coordinates $\tilde x_{ki}$ at reference ages
$\tilde a_{ki}$, non-centred speeds $z_{ki}$ with $\delta_{ki}=u_i^\top\psi_k+\sigma_{\delta k}z_{ki}$,
and flat levels $b_{ki}$. The departure age $\alpha_{ki}$ is a deterministic
function of these (\S\ref{sec:tau-map}).

\section{The data file}\label{sec:data}

\code{data/build\_joint\_dataset.R} reads the two source CSVs and writes
\code{data/adni\_joint.rds}, a list that is a strict superset of the original
\code{adni\_amyloid.rds}. The amyloid fields are rebuilt from \code{dta\_amy.csv}
and asserted equal to the original file to $10^{-10}$ (the test suite repeats this
assertion). The cohort is the 1{,}101 subjects with at least two amyloid scans, in
the original row order. Tau scans are joined on the ADNI \code{RID}; the column
\code{RID\_num} is numbered separately in each CSV and must not be used as a key.

\begin{center}
\begin{tabular}{@{}llp{0.55\linewidth}@{}}
\toprule
Field & Dimension & Definition \\
\midrule
\code{n\_obs\_tau} & $N$ & $n_{Ti}$; zero for the 453 subjects without tau scans \\
\code{y\_tau\_mat} & $N\times 7$ & $y_{Tij}$, zero-padded like \code{y\_mat} \\
\code{age\_tau\_mat} & $N\times 7$ & $a_{Tij}$, zero-padded \\
\code{a\_tilde\_tau} & $N$ & $\tilde a_{Ti}$ = mean tau scan age; for subjects without tau, the amyloid centroid $\tilde a_{Ai}$ (the proposal's ``fixed reference age'') \\
\code{x\_init\_tau} & $N$ & mean tau level, a starting value for $\tilde x_{Ti}$; the cohort median for subjects without tau \\
\code{has\_tau} & $N$ & logical \\
\code{n\_subj\_tau} & 1 & 648 \\
\code{rid} & $N$ & ADNI RID \\
\code{centiloid} & 2 & intercept and slope of SUVR $= a + b\,$CL; $c_A$ is this map at 24 CL \\
\bottomrule
\end{tabular}
\end{center}

Covariate centring is unchanged: APOE4 and sex indicators are centred over all
1{,}900 amyloid subjects in the source file, not over the cohort, so the shipped
amyloid posterior's intercepts remain on the same scale. Totals: 3{,}393 amyloid
scans, 1{,}452 tau scans on 648 subjects. No tau positivity threshold is stored in
the data; it is a setting (\S\ref{sec:settings}).

\code{subset\_subjects(dat, k)} (in \code{run\_model.R}) cuts \emph{every}
per-subject field (any vector of length $N$ or matrix with $N$ rows) to the rows
\code{k}, so a subsample stays self-consistent across amyloid and tau fields.

\section{Settings}\label{sec:settings}

\code{SETTINGS\_TAU} is \code{SETTINGS} with the entries below overridden, so
\code{make\_rate\_basis(SETTINGS\_TAU)} and \code{make\_var\_basis(y, SETTINGS\_TAU)}
build tau bases with the existing basis code. \code{SETTINGS\_JOINT} holds what is
shared. Tau is the NEO-T composite (volume-weighted flortaucipir SUVR of
bilateral inferior and middle temporal gyri; Ossenkoppele et al.\ 2022, Nat Med),
inferior cerebellar reference, no partial-volume correction. The tau entries are
calibrated from the literature on that measurement, not from these data: Jack et
al.\ 2020 (JAMA Neurol; amyloid-negative inferior temporal SUVR $1.18$, SD $0.08$;
rates $0.005$/yr below and $0.03$/yr above 68 CL), Ossenkoppele et al.\ 2022
(positivity = mean $+2$ SD of amyloid-negatives; 14.5\% of A+ unimpaired NEO-T+),
Ossenkoppele et al.\ 2025 (Nat Neurosci; GMM cut-off 1.40), Devous et al.\ 2018
and Timmers et al.\ 2020 (test--retest 2--4\%). Only \code{d\_coup} is marked P
(a modelling choice).

\begin{center}
\begin{tabular}{@{}llp{0.6\linewidth}@{}}
\toprule
Setting & Value & Role \\
\midrule
\code{grid\_lo}, \code{grid\_hi}, \code{grid\_step} & 0.70, 4.50, 0.002 & tau level grid (1{,}901 points); observed NEO-T range is 0.92--4.01 \\
\code{thresh} ($c_T$) & 1.34 & $1.18+2\times0.08$; anchors $G_T$ and truncates $b_{Ti}$, $\mu_{bT}$ \\
\code{theta1}, \code{theta1\_sd} & $-5.3$, 2.0 & prior mean/sd of $\theta_{T1}$: $e^{-5.3}=0.005$ SUVR/yr, the slow regime \\
\code{init\_rise} & 6 & initial rate curve rises $0.005\to0.03$/yr from floor to threshold (amyloid: 20) \\
\code{nu1\_mean} & $-3.5$ & $\log\omega_T$ at the low end; $e^{-3.5}=0.030$ SUVR, i.e.\ SD 0.043 per scan \\
\code{mu\_b\_mean}, \code{mu\_b\_sd} & 1.17, 0.05 & prior on $\mu_{bT}$, truncated to $(\code{grid\_lo}, c_T)$ \\
\code{sigma\_b\_sd}, \code{sigma\_b\_max} & 0.10, 0.30 & half-normal scale and cap for $\sigma_{bT}$ (latent spread $\approx0.07$) \\
\code{phi1\_mean}, \code{phi1\_sd} & 0.5, 1 & conditional tau susceptibility intercept: $q=0.62$, 95\% range 0.18--0.93 \\
\code{g1\_age} & 75 & centre of the tau AFT intercept, $\gamma_{T1}\sim N(\log(75-25),0.75^2)$ (amyloid: 70) \\
\midrule
\code{d\_coup} & 0.10 (P) & sigmoid width of $H_A$, in amyloid SUVR \\
\code{beta\_sd} & 1.0 & $\beta\sim N(0,1)$ truncated to $\beta\ge 0$ \\
\code{rho\_beta} & (2, 2) & $(\rho+1)/2\sim\mathrm{Beta}(2,2)$ \\
\code{subj\_block} & \code{"four"} & per-subject sampler blocking (\S\ref{sec:samplers}) \\
\bottomrule
\end{tabular}
\end{center}

Everything not listed (age bounds 25--125, $r_{\min}$, $t$ degrees of freedom 4,
random-walk truncations, $z$ bound 4, AFT intercept prior centred at age 70)
is inherited from \code{SETTINGS} for both biomarkers.

\section{The coupling and the exposure clock}\label{sec:exposure}

\subsection{The multiplier}

The proposal couples tau to amyloid through a rate multiplier
\begin{equation}
m(y)=\exp\{s(y)\},\qquad s(y)=\beta\,[H_A(y)-H_A(c_A)],\qquad
H_A(y)=\expit\!\left(\frac{y-c_A}{d}\right),
\label{eq:m}
\end{equation}
with $d=\code{d\_coup}$. Since $H_A(c_A)=\tfrac12$, $m(c_A)=1$ and
$m\in(e^{-\beta/2},e^{\beta/2})$. Implemented as \code{nfCouple(y, beta, cA, d)}.
The tau ODE for an accumulator is
\begin{equation}
\frac{dT_i}{da}=e^{\delta_{Ti}}\,m\!\big(A_i(a)\big)\,\rbar_T\!\big(T_i(a)\big),\qquad T_i(\alpha_{Ti})=\mu_{bT},
\label{eq:tauode}
\end{equation}
where $A_i(a)$ is the subject's \emph{candidate accumulator} amyloid path
(flat at $\mu_{bA}$ before $\alpha_{Ai}$, then on the amyloid clock). Dividing by
$\rbar_T$ and integrating gives
\begin{equation}
G_T\!\big(T_i(a)\big)=G_T(\mu_{bT})+e^{\delta_{Ti}}\int_{\alpha_{Ti}}^{a} m\!\big(A_i(u)\big)\,du .
\label{eq:tausol}
\end{equation}

\subsection{The exposure clock: the one deviation from the proposal}

The proposal evaluates the integral in \eqref{eq:tausol} by per-subject age
quadrature. The implementation does not. Define the exposure accumulated since
\emph{amyloid} departure,
\[
\Psi_i(u)=\int_{\alpha_{Ai}}^{u} m\!\big(A_i(v)\big)\,dv ,
\]
and note that $A_i$ is constant ($=\mu_{bA}$) before $\alpha_{Ai}$ and, after it,
satisfies $dA_i=e^{\delta_{Ai}}\rbar_A(A_i)\,dv$, so $dv=e^{-\delta_{Ai}}dA_i/\rbar_A(A_i)$.
Hence
\begin{equation}
\Psi_i(u)=\begin{cases}
m(\mu_{bA})\,(u-\alpha_{Ai}), & u\le\alpha_{Ai},\\[3pt]
e^{-\delta_{Ai}}\big[\E(A_i(u))-\E(\mu_{bA})\big], & u>\alpha_{Ai},
\end{cases}
\qquad
\E(y)=\int_{\ell_A}^{y}\frac{m(v)}{\rbar_A(v)}\,dv .
\label{eq:psi}
\end{equation}
$\E$ is a \emph{population} table on the amyloid level grid: it depends on
$\theta_A$ and $\beta$ only, and is tabulated once per iteration by the
cumulative trapezoid rule, exactly as $G_A$ is. The lower limit $\ell_A$ is the
grid bottom; only differences of $\E$ ever appear. A subject's exposure between
two ages is then two lookups, and its inverse is a monotone table inverse of
$\E$ followed by the amyloid clock:
\begin{equation}
\Psi_i^{-1}(\psi)=\begin{cases}
\alpha_{Ai}+\psi/m(\mu_{bA}), & \psi\le0,\\[3pt]
\alpha_{Ai}+e^{-\delta_{Ai}}\big[G_A(y^\ast)-G_A(\mu_{bA})\big],\quad
y^\ast=\E^{-1}\!\big(\E(\mu_{bA})+e^{\delta_{Ai}}\psi\big), & \psi>0 .
\end{cases}
\label{eq:psiinv}
\end{equation}
The proposal's $C_i(a)=\int_{\tilde a_{Ti}}^{a}m_i(u)\,du$ is
$\Psi_i(a)-\Psi_i(\tilde a_{Ti})$, and
\eqref{eq:tausol} becomes $G_T(T_i(a))=G_T(\mu_{bT})+e^{\delta_{Ti}}[\Psi_i(a)-\Psi_i(\alpha_{Ti})]$.
At $\beta=0$, $m\equiv1$, $\E=G_A+\text{const}$, and both cases of \eqref{eq:psi}
reduce to $\Psi_i(u)=u-\alpha_{Ai}$: the uncoupled model is recovered exactly, which
the tests verify to $5\times10^{-13}$.

\paragraph{Functions.}
\begin{itemize}[nosep]
\item \code{nfExpClock(r, step, n, lo, r\_min, beta, cA, d)}: returns $\E$ on the grid.
\item \code{nfAmyState(xA, deltaA, atilA, mu\_bA, G, rgrid, lo, step, ngrid, r\_min)}:
returns $(g_0, G_b, \alpha_{Ai})$ with $g_0=G_{A,\mathrm{ext}}(\tilde x_{Ai})$,
$G_b=G_A(\mu_{bA})$, $\alpha_{Ai}=\tilde a_{Ai}-e^{-\delta_{Ai}}(g_0-G_b)$; returns
$\alpha_{Ai}=-\infty$ if $\tilde x_{Ai}$ is above the table.
\item \code{nfExposure(u, g0, Gb, alphaA, deltaA, atilA, mu\_bA, G, E, lo, step, ngrid, beta, cA, d)}:
$\Psi_i(u)$ by \eqref{eq:psi}; \code{NaN} if $A_i(u)$ is off the table.
\item \code{nfInvExposure(psi, g0, Gb, alphaA, deltaA, mu\_bA, G, E, lo, step, ngrid, beta, cA, d)}:
returns $(\Psi_i^{-1}(\psi),\,A_i(\Psi_i^{-1}(\psi)))$ by \eqref{eq:psiinv}, the
second component so the caller gets $m_i$ at that age without another lookup.
\end{itemize}

\section{The tau location coordinate and its density}\label{sec:tau-map}

\subsection{The map from location to departure age}

For amyloid, the existing map is $\alpha_{Ai}=\tilde a_{Ai}-e^{-\delta_{Ai}}\Delta G_{Ai}$,
$\Delta G_{Ai}=G_{A,\mathrm{ext}}(\tilde x_{Ai})-G_A(\mu_{bA})$, with Jacobian
$J_{Ai}=e^{-\delta_{Ai}}/\rbar_A(\cdot)$ (evaluated at $\tilde x_{Ai}$ above the
floor, at $\mu_{bA}$ below it). This is \code{dxtilde} and is unchanged.

For tau, setting $a=\tilde a_{Ti}$ in the solution and using $G_{T,\mathrm{ext}}(\tilde x_{Ti})=G_T(T_i(\tilde a_{Ti}))$,
\begin{equation}
\Psi_i(\alpha_{Ti})=\Psi_i(\tilde a_{Ti})-e^{-\delta_{Ti}}\Delta G_{Ti},\qquad
\Delta G_{Ti}=G_{T,\mathrm{ext}}(\tilde x_{Ti})-G_T(\mu_{bT}),
\label{eq:taumap}
\end{equation}
so $\alpha_{Ti}=\Psi_i^{-1}\big(\Psi_i(\tilde a_{Ti})-e^{-\delta_{Ti}}\Delta G_{Ti}\big)$.
Differentiating \eqref{eq:taumap} in $\tilde x_{Ti}$ with the amyloid path and both
speeds fixed, $m(A_i(\alpha_{Ti}))\,\partial\alpha_{Ti}/\partial\tilde x_{Ti}=-e^{-\delta_{Ti}}D_T(\tilde x_{Ti})$, giving
\begin{equation}
J_{Ti}=\frac{e^{-\delta_{Ti}}\,D_T(\tilde x_{Ti})}{m\big(A_i(\alpha_{Ti})\big)},\qquad
D_T(x)=\begin{cases}1/\rbar_T(x), & x\ge\mu_{bT},\\ 1/\rbar_T(\mu_{bT}), & x<\mu_{bT}.\end{cases}
\label{eq:jac}
\end{equation}
The factor $1/m(A_i(\alpha_{Ti}))$ is the ``required change'' of the proposal's
\S7. The transformation $(\tilde x_{Ai},\tilde x_{Ti})\mapsto(\alpha_{Ai},\alpha_{Ti})$
is triangular ($\alpha_{Ai}$ does not depend on $\tilde x_{Ti}$), so the determinant
is $J_{Ai}J_{Ti}$ and the cross-derivative, though nonzero, does not enter.
As in \code{dxtilde}, $D$ uses the interpolated $1/\rbar$ at the point rather
than the slope of the piecewise-linear clock table; the two differ at the level of
the trapezoid discretisation.

\subsection{Departure-age marginals and the copula}

Both departure ages have the existing shifted, upper-truncated Weibull AFT
marginal with $\kappa_k=1/\sigma_{\alpha k}$, $\lambda_{ki}=\exp(x_i^\top\gamma_k)$,
support $(a_{\min},a_{\max}]=(25,125]$:
\[
f_{ki}(a)=\frac{\kappa_k}{\lambda_{ki}C_{ki}}\Big(\tfrac{a-a_{\min}}{\lambda_{ki}}\Big)^{\kappa_k-1}
\exp\Big[-\Big(\tfrac{a-a_{\min}}{\lambda_{ki}}\Big)^{\kappa_k}\Big],\quad
F_{ki}(a)=\frac{1-\exp[-((a-a_{\min})/\lambda_{ki})^{\kappa_k}]}{C_{ki}},
\]
$C_{ki}=1-\exp[-((a_{\max}-a_{\min})/\lambda_{ki})^{\kappa_k}]$.
Implemented as \code{nfLogTWeib} (log density) and \code{nfTWeibCDF} (CDF), both
using \code{nfExpm1} (a cancellation-free $e^x-1$, since \code{expm1} is not in
the NIMBLE DSL). The joint departure density is
\[
f_{\alpha i}(a_A,a_T)=c_\rho\big(F_{Ai}(a_A),F_{Ti}(a_T)\big)\,f_{Ai}(a_A)\,f_{Ti}(a_T),
\]
with the Gaussian copula density, $v_k=\Phi^{-1}(p_k)$,
\[
\log c_\rho(p_1,p_2)=-\tfrac12\log(1-\rho^2)-\frac{\rho^2(v_1^2+v_2^2)-2\rho v_1v_2}{2(1-\rho^2)},
\]
implemented as \code{nfLogCopula(p1, p2, rho)}. There is no speed copula; $z_{Ai}$
and $z_{Ti}$ are independent truncated standard normals.

\subsection{\code{dxtildeT}: the density of $\tilde x_{Ti}$}

The node $\tilde x_{Ti}\sim\code{dxtildeT}(\cdots)$ has log density
\begin{equation}
\log f_{Ti}(\alpha_{Ti})+\log c_\rho\big(F_{Ai}(\alpha_{Ai}),F_{Ti}(\alpha_{Ti})\big)
-\delta_{Ti}+\log D_T(\tilde x_{Ti})-\log m\big(A_i(\alpha_{Ti})\big),
\label{eq:dxt}
\end{equation}
evaluated in the order: $\Delta G_{Ti}$ and $\log D_T$ from the tau clock;
$(g_0,G_b,\alpha_{Ai})$ from \code{nfAmyState}; $\Psi_i(\tilde a_{Ti})$ from
\code{nfExposure}; $(\alpha_{Ti},A_i(\alpha_{Ti}))$ from \code{nfInvExposure};
then the three terms. It returns $-\infty$ if $\tilde x_{Ti}$ or $\tilde x_{Ai}$ is
above its table, if either implied age is outside $(a_{\min},a_{\max}]$, if either
CDF is exactly 0 or 1, or if any amyloid lookup leaves the table. The amyloid
marginal $f_{Ai}(\alpha_{Ai})J_{Ai}$ is supplied by \code{dxtilde} on
$\tilde x_{Ai}$, so the copula term appears exactly once in the posterior. No other
prior is placed on $\tilde x_{Ti}$.

Argument list (all scalars unless noted): tau --- \code{deltaT, atilT, mu\_bT, xgT,
sigma\_alphaT, aminT, amaxT, GT[], rT[], loT, stepT, ngridT, r\_minT}; amyloid ---
\code{xA, deltaA, atilA, mu\_bA, xgA, sigma\_alphaA, aminA, amaxA, GA[], rA[], EA[],
loA, stepA, ngridA, r\_minA}; coupling --- \code{beta, cA, d\_coup, rho}.
Here \code{xgk}$=x_i^\top\gamma_k$. The simulation function \code{rxtildeT} is a
stub (NIMBLE requires one; the model never simulates this node).

\section{The three-state subject likelihood: \code{dsubjJoint}}\label{sec:densities}

The data node for subject $i$ is the concatenated vector
$y_{i\cdot}=(y_{Ai1},\dots,y_{Ai,\max n_A},\,y_{Ti1},\dots,y_{Ti,\max n_T})$,
zero-padded, with the tau block starting at offset \code{offT}$=\max n_A$.
Writing $t_4(y;\mu,\omega)$ for the scaled Student-$t$ log density
(\code{nfLogDT}; $\omega$ is a scale, not an SD), the four branch sums are
\begin{align*}
\ell^{0}_{Ai}&=\sum_j t_4\big(y_{Aij};b_{Ai},\omega_A(b_{Ai})\big), &
\ell^{1}_{Ai}&=\sum_j t_4\big(y_{Aij};A_i(a_{Aij}),\omega_A(A_i(a_{Aij}))\big),\\
\ell^{0}_{Ti}&=\sum_j t_4\big(y_{Tij};b_{Ti},\omega_T(b_{Ti})\big), &
\ell^{1}_{Ti}&=\sum_j t_4\big(y_{Tij};T_i(a_{Tij}),\omega_T(T_i(a_{Tij}))\big),
\end{align*}
with $A_i(a)=G_A^{-1}\big(\max\{G_b,\;g_0+e^{\delta_{Ai}}(a-\tilde a_{Ai})\}\big)$
as in \code{dsubjMix}, and
\begin{equation}
T_i(a)=G_T^{-1}\Big(\max\big\{G_T(\mu_{bT}),\;G_{T,\mathrm{ext}}(\tilde x_{Ti})+e^{\delta_{Ti}}[\Psi_i(a)-\Psi_i(\tilde a_{Ti})]\big\}\Big),
\label{eq:Tpath}
\end{equation}
which is the proposal's eq.~(21) in exposure-clock form: the inner $\max$ holds
tau at its floor before tau departure. With $\pi_i=\expit(w_i^\top\phi_A)$ and
$q_i=\expit(w_i^\top\phi_T)$,
\begin{equation}
\log L_i=\logsumexp\Big\{
\underbrace{\log(1-\pi_i)+\ell^0_{Ai}+\ell^0_{Ti}}_{\ell_{00}},\;
\underbrace{\log\pi_i+\log(1-q_i)+\ell^1_{Ai}+\ell^0_{Ti}}_{\ell_{10}},\;
\underbrace{\log\pi_i+\log q_i+\ell^1_{Ai}+\ell^1_{Ti}}_{\ell_{11}}\Big\}.
\label{eq:lik}
\end{equation}
State 11 uses the \emph{same} amyloid accumulator path in $\ell^1_{Ti}$ as in
$\ell^1_{Ai}$; there is no tau-only state. When $n_{Ti}=0$ the tau sums are zero and
\eqref{eq:lik} collapses to the amyloid two-branch mixture (verified numerically to
$10^{-10}$ against \code{dsubjMix}). The function returns $-\infty$ if either
location is at or above its table top, or if any trajectory lookup leaves a table.
Argument order: \code{nA, agesA[], nT, agesT[], offT}; the amyloid block
(\code{xA, atilA, deltaA, bA, logit\_pi, mu\_bA, GA[], rA[], EA[], loA, stepA, ngridA,
r\_minA, omA[], vloA, vstepA, nvA}); the tau block (same pattern with \code{xT, atilT,
deltaT, bT, logit\_q, mu\_bT, GT[], rT[], loT, stepT, ngridT, r\_minT, omT[], vloT,
vstepT, nvT}); then \code{beta, cA, d\_coup, df}.

\section{The model: \code{build\_joint\_model()}}\label{sec:model}

\code{build\_joint\_model(dat, bA, vA, bT, vT, SA, ST, SJ)} returns
\code{list(code, constants, data, inits)} for \code{nimbleModel}. The amyloid
population block is copied verbatim from \code{R/04\_model.R}. New nodes:

\paragraph{Amyloid, one addition.}
\code{E\_grid[1:ngrid] <- nfExpClock(r\_grid, y\_step, ngrid, y\_lo, r\_min, beta, thresh, d\_coup)}.

\paragraph{Tau population.} The same structure as amyloid with suffix \code{\_T}:
\begin{itemize}[nosep]
\item Rate curve: $\theta_{T1}\sim N(\code{theta1\_T},\code{theta1\_sd\_T}^2)$ \textbf{free}
(the amyloid $\theta_1$ is pinned because its basis support lies below the floor;
the tau basis function 1 has support $[0.70,1.10]$, which overlaps the plausible
$\mu_{bT}$, so pinning it would be an accidental prior). Then
$\theta_{Tj}\mid\theta_{T,j-1}\sim N(\theta_{T,j-1},\lambda_{Tj}/\varsigma_T)$, $j=2,\dots,10$,
with the same truncated $\mathrm{Exp}(1)$ local and $\mathrm{Gamma}(1,1)$ global
priors (\code{lam\_sp\_th\_T}, \code{sp\_th\_T}). \code{r\_grid\_T}, \code{G\_grid\_T} as for amyloid.
\item Measurement scale: $\nu_{T1}\sim N(\code{nu1\_mean\_T},1)$, fixed-width random walk
(sd 0.5), all truncated to $(-9,1)$. \code{omega\_grid\_T}.
\item Floor: $\mu_{bT}\sim N(1.05,0.10^2)$ truncated to $(\code{y\_lo\_T},c_T)$.
\item Departure: $\gamma_{T1}\sim N(\log 45,0.75^2)$, $\gamma_{Tq}\sim N(0,0.25^2)$,
$\sigma_{\alpha T}\sim\mathrm{Gamma}(2,8)$ truncated to $(0,2)$.
\item Speed: $\psi_{Tq}\sim N(0,0.25^2)$ (no intercept), $\sigma_{\delta T}\sim N^+(0,1)$ truncated at 2.
\item Conditional susceptibility: $\phi_{T1}\sim N(0,1)$, $\phi_{Tq}\sim N(0,0.5^2)$;
\code{q\_bar} $=\expit(\phi_{T1})$.
\item $\sigma_{bT}\sim N^+(0,0.15)$ truncated at 0.5.
\end{itemize}

\paragraph{Coupling and dependence.}
$\beta\sim N(0,\code{beta\_sd}^2)$ truncated to $\beta\ge0$;
$u_\rho\sim\mathrm{Beta}(2,2)$, \code{rho\_alpha} $=2u_\rho-1$.

\paragraph{Per subject.} In addition to the amyloid nodes
\code{z[i], delta[i], xg[i], lp\_i[i], x\_tilde[i], b[i]}:
\begin{lstlisting}
z_T[i]     ~ T(dnorm(0, 1), -z_bound, z_bound)
delta_T[i] <- inprod(X[i,], psi_T) + sigma_delta_T * z_T[i]
xg_T[i]    <- inprod(Xg[i,], gamma_T)          # AFT log scale
lq_i[i]    <- inprod(Xg[i,], phi_T)            # logit q_i
x_tilde_T[i] ~ dxtildeT(...)                   # eq. (8)
b_T[i]     ~ T(dnorm(mu_b_T, sigma_b_T), , thresh_T)
y_joint[i, 1:max_nn] ~ dsubjJoint(...)         # eq. (10)
\end{lstlisting}

\paragraph{Constants and data.} All amyloid constants as before, plus
\code{n\_obs\_T, max\_n\_T, max\_nn}$=\max n_A+\max n_T$, \code{age\_mat\_T,
a\_tilde\_T}, the tau bases (\code{K\_T, ngrid\_T, Bh\_T, y\_lo\_T, y\_step\_T, Kv\_T,
nvgrid\_T, Bv\_T, v\_lo\_T, v\_step\_T}), tau prior hyperparameters, and
\code{d\_coup, beta\_sd, rho\_a, rho\_b}. Data: \code{y\_joint = cbind(y\_mat,
y\_tau\_mat)} and the pinned amyloid $\theta_1$.

\paragraph{Initial values.} Amyloid inits as before. Tau: $\theta_T$ from a
least-squares fit to a log-rate target that rises by a factor 20 from the grid
bottom to $c_T$ and then flattens (the same rule as amyloid); $\nu_T$ at
\code{nu1\_mean\_T}; $\mu_{bT}=\min(1.05,c_T-0.05)$; $\gamma_{T1}=\log 45$;
$\sigma_{\alpha T}=0.27$; $\sigma_{\delta T}=0.4$; $\phi_T=0$; $\sigma_{bT}=0.08$;
$\tilde x_{Ti}$ = \code{x\_init\_tau} clipped into the grid; $b_{Ti}=\min(\code{x\_init\_tau},c_T-0.01)$;
$\beta=0$ and $u_\rho=0.5$ (uncoupled, independent start). Speeds $z_{Ti}$ are set
by the same feasibility rule as $z_{Ai}$ (minimum speed so that the implied
departure age exceeds $a_{\min}$), which is exact at $\beta=0$ because the tau map
is then the plain clock map. \code{feasible\_z()} is the shared helper.

\section{Runner, samplers and post-processing}\label{sec:runner}

\subsection{\code{configure\_joint\_samplers()}}\label{sec:samplers}

Monitors: every population parameter of both biomarkers plus \code{beta,
rho\_alpha, pi\_bar, q\_bar}; \code{monitors2} (thinned 20$\times$):
\code{x\_tilde, z, b, x\_tilde\_T, z\_T, b\_T}. Default samplers are removed before
any replacement is added (the ordering issue documented in \code{run\_model.R}).
Blocks:
\begin{itemize}[nosep]
\item \code{AF\_slice}: \code{nu}; \{\code{gamma[1]}, \code{theta[2:10]}\}; \code{gamma[2:3]}; \code{psi};
\code{nu\_T}; \{\code{gamma\_T[1]}, \code{theta\_T[1:10]}\}; \code{gamma\_T[2:3]}; \code{psi\_T}; \code{phi\_T}.
\item \code{slice}: \code{sigma\_delta, sigma\_alpha, mu\_b, sigma\_b} and their
\code{\_T} twins, \code{beta}, \code{u\_rho}; per subject \code{b[i]}, \code{b\_T[i]}.
\item Per subject: with \code{subj\_block = "four"} (default), one \code{RW\_block}
over $(\tilde x_{Ai},z_{Ai},\tilde x_{Ti},z_{Ti})$ for subjects with tau scans and two
pairs $(\tilde x_{Ai},z_{Ai})$, $(\tilde x_{Ti},z_{Ti})$ otherwise; with \code{"pairs"},
two pairs for everyone. The four-block exists because an amyloid proposal moves the
implied tau departure age through the exposure clock.
\item \code{phi} (amyloid susceptibility) keeps NIMBLE's default, as in the amyloid runner.
\end{itemize}
An assertion fails if any node is left without a sampler.

\subsection{\code{fit\_joint\_model(tier, n\_chains, seed, out\_dir, verbose, jitter, n\_iter, chunk)}}

Tiers: \code{smoke} (250 subjects, 300 iterations), \code{struct} (all, 800),
\code{infer} (all, 20{,}000); \code{n\_iter} and \code{chunk} override. Builds both
bases per biomarker (\code{make\_var\_basis} on that biomarker's observed scans),
checks the log posterior at the inits is finite and names offending nodes if not,
compiles, and runs chains sequentially in chunks, writing
\code{<tier>\_chains.rds} and \code{<tier>\_subj.rds} after every chunk and
\code{<tier>\_fit.rds} at the end (a list with \code{chains, subj, dat, bA, vA, bT, vT,
settings, minutes}). Chains $2,\dots$ start from the inits jittered by
\code{jitter}$\cdot N(0,1)$ on $\theta,\nu,\psi,\gamma$ and their tau twins and $\phi_T$,
$0.02\,N(0,1)$ on both floors, $|0.2\,N(0,1)|$ on $\beta$ and $0.1\,N(0,1)$ on $u_\rho$,
with the jitter scale halved until the state is feasible.

\subsection{\code{joint\_state\_probs(fit, max\_draws)}}

Posterior state probabilities $r_{si}=P(S_i=s\mid\text{data},\text{draw})$, averaged
over up to \code{max\_draws} saved subject draws. At each draw it rebuilds
$r_A,G_A,\E,r_T,G_T,\omega_A,\omega_T$ in R from the population draw, recomputes
the subject's paths with \code{joint\_paths()} (the R twin of \code{nfExposure} and
\eqref{eq:Tpath}, using \code{sa\_interp}/\code{sa\_inv\_monotone}), evaluates the three
branches of \eqref{eq:lik} with \code{dt\_log()}, and normalises. Returns an
$N\times3$ matrix with columns \code{00, 10, 11}; rows sum to one. Biological
summaries for a state should be weighted by the corresponding column, as the
proposal's \S8 specifies.

\section{Tests: \code{tests/test\_joint.R}}\label{sec:tests}

Compiles the new functions standalone and checks, with a bent rate curve
(the inits plus a fixed perturbation) and a hand-built subject:
\begin{enumerate}[nosep]
\item \textbf{Uncoupled limit.} $\E-G_A$ constant at $\beta=0$; \code{dxtildeT} at
$\beta=0,\rho=0$ equals \code{dxtilde} called with tau arguments, at three
locations; the copula term vanishes at $\rho=0$.
\item \textbf{Constant-exposure limit.} A subject at 1.30 SUVR at age 70, moving
slowly, has $H_A>0.996$ over a tau window 72--80: the multiplier is constant to
1\%, the exposure grows linearly at that rate, and the rate exceeds $e^{0.45\beta}$.
\item \textbf{Round trips and Jacobian.} $\tilde x_T\to\alpha_T\to\tilde x_T$ at four
locations (including one below the floor) to $10^{-6}$; the Jacobian read off
\eqref{eq:dxt} against central finite differences of $\alpha_T(\tilde x_T)$ to 2\%;
a nonzero amyloid-to-tau cross-derivative.
\item \textbf{ODE agreement.} \eqref{eq:Tpath} against an RK4 solve of
\eqref{eq:tauode} with step 0.01 yr at three ages, to $3\times10^{-3}$ SUVR
(observed agreement: 4 decimals).
\item \textbf{Normalisation.} Truncated Weibull density integrates to one; its CDF is
one at $a_{\max}$; the copula density integrates to one over the unit square.
\item \textbf{Marginalisation.} For a subject with no tau scans, \code{dsubjJoint}
equals \code{dsubjMix} to $10^{-10}$; state weights sum to one.
\end{enumerate}
All 24 pass. \code{tests/test\_package.R} (the amyloid-only suite) passes 26 of 26
with the new files sourced.

\section{Status, and what is deliberately not done}

\begin{itemize}[nosep]
\item Compiled and smoke-run on 250 subjects, 2 chains $\times$ 300 iterations:
compile 0.9 min, about 0.12 s per iteration; all parameters move. On the full cohort
expect roughly 0.5 s per iteration. No inference is claimed from this run.
\item \textbf{No tail continuation.} A trajectory that leaves a level table returns
$-\infty$, as in the amyloid model; the proposal asks for constant extension of
$\log r$ and $\log\omega$ beyond the grid. The tau grid is wide (0.70--4.50) to make
this rare for the auxiliary tau coordinates of the 454 subjects without tau scans,
but it remains a parameter-dependent support restriction.
\item \textbf{Tau summary.} The model was first run on the mesial temporal
meta-ROI, for which the three-state restriction $Z_T\le Z_A$ is biologically wrong
(age-related medial temporal tau without amyloid): that fit forced tau-positive,
amyloid-flat subjects into the latent-amyloid state and pulled the amyloid
susceptibility $\bar\pi$ from 0.74 to 0.91. The main dataset is now the NEO-T
neocortical composite, for which the restriction is defensible, with priors
recalibrated from the literature (\S\ref{sec:settings}). Only $d_{\mathrm{coup}}$
remains a provisional choice.
\item Chains now run in parallel (one process each). Not implemented: a full
recovery study (single-chain checks at $\beta\in\{0,1.75\}$ exist), model
comparison (a)--(d), prediction hold-outs, prior sensitivity, a flexible coupling
spline, and departure ordering.
\end{itemize}

\end{document}
